% Part: second-order-logic
% Chapter: sol-and-sets
% Section: comparing-sets

\documentclass[../../../include/open-logic-section]{subfiles}

\begin{document}

\olfileid[fr]{sol}{set}{cmp}

\olsection{Comparer les ensembles}

\begin{prop}
La !!{formula}~$\lforall[x][(X(x)\lif Y(x))]$
définit l'inclusion :
$\Sat{M}{\lforall[x][(X(x)\lif Y(x))]}[s]$
si et seulement si $s(X)\subseteq s(Y)$.
\end{prop}

\begin{prop}
La !!{formula}~$\lforall[x][(X(x)\liff Y(x))]$
définit l'égalité des ensembles :
$\Sat{M}{\lforall[x][(X(x)\liff Y(x))]}[s]$
si et seulement si $s(X)=s(Y)$.
\end{prop}

\begin{prop}
La !!{formula}~$\lexists[x][X(x)]$ définit
la propriété d'être non vide :
$\Sat{M}{\lexists[x][X(x)]}[s]$
si et seulement si $s(X)\neq\emptyset$.
\end{prop}

Un ensemble~$X$ n'est pas plus grand qu'un
ensemble~$Y$, ce que l'on note
$\cardle{X}{Y}$, s'il existe une fonction
!!{injective} $f\colon X\to Y$.
La logique du second ordre peut exprimer que
la restriction d'une fonction au sous-ensemble
$X$ est !!{injective} et que ses valeurs sur~$X$
appartiennent à~$Y$. Elle peut donc définir cette
relation entre sous-ensembles du !!{domain}.

\begin{prop}
La !!{formula}
\begin{multline*}
\lexists[u][(\lforall[x][(X(x)\lif Y(u(x)))]\land{}\\
\lforall[x][\lforall[y][((X(x)\land X(y)\land
\eq[u(x)][u(y)])\lif\eq[x][y])]])]
\end{multline*}
définit la relation $\cardle{X}{Y}$.
\end{prop}

Deux ensembles ont la même taille, ou sont
\emph{équipotents}, ce que l'on note
$\cardeq{X}{Y}$, s'il existe une
!!{bijection} $f\colon X\to Y$.

\begin{prop}
La !!{formula}
\begin{multline*}
\lexists[u][(\lforall[x][(X(x)\lif Y(u(x)))]\land{}\\
\lforall[x][\lforall[y][((X(x)\land X(y)\land
\eq[u(x)][u(y)])\lif\eq[x][y])]]\land{}\\
\lforall[y][(Y(y)\lif\lexists[x][(X(x)
\land\eq[y][u(x)])])])]
\end{multline*}
définit la relation $\cardeq{X}{Y}$.
\end{prop}

Nous abrégerons désormais ces !!{formula}s,
dans cet ordre, par $X\subseteq Y$, $X=Y$,
$X\neq\emptyset$, $\cardle{X}{Y}$ et
$\cardeq{X}{Y}$. Cette notation peut prêter
à confusion, puisque nous l'employons aussi
pour parler informellement des ensembles~$X$
et~$Y$. Ici, elle abrège des !!{formula}s du
second ordre contenant les !!{variable}s de
relation unaires $X$ et~$Y$.

\begin{prop}
L'!!{sentence}
$\lforall[X][\lforall[Y][((\cardle{X}{Y}\land
\cardle{Y}{X})\lif\cardeq{X}{Y})]]$
est valide.
\end{prop}

\begin{proof}
Dans toute !!{structure}~$\Struct{M}$,
pour tous $X,Y\subseteq\Domain{M}$,
si $\cardle{X}{Y}$ et $\cardle{Y}{X}$,
alors $\cardeq{X}{Y}$. C'est le théorème
de Schröder--Bernstein, valable pour tous
ensembles~$X$ et~$Y$.
\end{proof}

L'injectivité requise dans ces deux formules porte sur les
arguments de~$X$. Toute fonction de $X$ dans~$Y$ se prolonge
au !!{domain} entier en donnant à tous les arguments hors
de~$X$ une même valeur de ce domaine, qui est non vide.
Le prolongement n'a pas à être injectif hors de~$X$.

Par exemple, pour $\Domain{M}=\Nat$, les ensembles
$X=\Nat\setminus\{0\}$ et $Y=\Nat$ sont équipotents.
Une bijection de $X$ sur~$Y$ utilise déjà tous les éléments de
$\Nat$ comme valeurs sur~$X$. Si son prolongement devait être
injectif sur tout le domaine, aucune valeur ne serait disponible
pour l'argument~$0$.

\end{document}
